% !PRINTED 171
\weberpage{171}

%<<CONT
also:
\[
\mathrm{Lim}\Bigl(\log\frac{\varkappa^2}{16} - i\pi\omega\Bigr) = 0,
\]
andererseits ist $2iK' = \pi\vartheta_{00}^2\,\omega$ (\S\,42) oder:
\[
2K' + \log\frac{\varkappa^2}{16} = \Bigl(\log\frac{\varkappa^2}{16} - i\pi\omega\Bigr) + i\pi\omega(1-\vartheta_{00}^2),
\]
woraus:
\weq{15}{\mathrm{Lim}_{\varkappa^2=0}\Bigl(K' + \frac12\log\frac{\varkappa^2}{16}\Bigr) = 0.}
Da nun $K$ und $K'$ partikulare Integrale von (8) sind (für
$x=\varkappa^2$), so haben beide die Form:
\[
K = a_1 y_1 + a_2 y_2, \qquad K' = a_1' y_1 + a_2' y_2,
\]
worin $a_1, a_2, a_1', a_2'$ konstant sind. Da $K$ für $x=0$ endlich
bleibt und den Wert $\frac12\pi$ erhält, so ist $a_2=0$, $a_1=\frac12\pi$, und
damit $K'+\frac12\log x$ endlich bleibe und [nach (15)] den Wert $\log 4$
erhalte, muß $a_2'=1$ und $a_1'=0$ sein. Also ergibt sich, wenn
man $x=\varkappa^2$, $1-x=\varkappa'^2$ setzt:
\weq{16}{K = \frac{\pi}{2}F(\varkappa^2), \qquad K' = \frac12 G(\varkappa^2) - \frac12\log\frac{\varkappa^2}{16\varkappa'^2}F(\varkappa^2)}
und
\weq{17}{q = e^{-\frac{\pi K'}{K}} = \frac{\varkappa^2}{16\varkappa'^2}\,e^{-\frac{G(\varkappa^2)}{F(\varkappa^2)}},}
worin, wie schon bemerkt, die Reihen $G(\varkappa^2)$, $F(\varkappa^2)$ konvergent
sind, so lange der absolute Wert von $\varkappa^2$ ein echter Bruch ist.
In
der Differentialgleichung (4) liegen nun freilich die Mittel, die
gefundenen Ausdrücke für $K,K',q$ über dies Konvergenzbereich % [sic] The source prints “dies” (for “diesen”).
hinaus fortzusetzen. Einfacher gelangt man dazu aber durch die
Anwendung der \gesperrt{Landen}schen \gesperrt{Transformation}.

Wir begrenzen die Ebene der komplexen Werte $\varkappa^2$, indem
wir längs der reellen Achse zwei Schnitte von 0 bis $-\infty$ und
von 1 bis $\infty$ legen. In der so begrenzten Ebene sind dann alle
Wurzeln aus $\varkappa^2, \varkappa'^2$ eindeutig dadurch bestimmt, daß sie für
positive echt gebrochene $\varkappa^2$ \gesperrt{reell und positiv} sein sollen. Es
ist nun nach den Formeln der \gesperrt{Landen}schen Transformation
[\S\,44, (28)], wenn wir mit $\varkappa_1,\varkappa_1',K_1,K_1'$ die Funktionen von $\omega$
bezeichnen, die sich aus $\varkappa,\varkappa',K,K'$ ergeben, wenn $\omega$ durch $2\omega$
ersetzt wird:
\weq{18}{\begin{aligned}
\varkappa_1 &= \frac{1-\varkappa'}{1+\varkappa'}, \qquad \varkappa_1' = \frac{2\sqrt{\varkappa'}}{1+\varkappa'}, \\
2K_1 &= (1+\varkappa')K, \qquad K_1' = (1+\varkappa')K'.
\end{aligned}}
